\documentclass{beamer}

\usepackage[utf8]{inputenc}
\usepackage[english]{babel}
\usepackage{textpos}
\usepackage{graphicx}
\usepackage{textcomp}
\usepackage{animate}
\usepackage{comment}
%\usepackage{multicol}
%\setlength{\columnsep}{5cm}
\usepackage{amsmath}
\usepackage{amsfonts}
\usepackage{verbatim}
\usepackage{varwidth}
\usepackage{caption}
\usepackage{dcolumn}
\usepackage{amssymb}
\usepackage{algorithm}
\usepackage{algorithmic}
%\usepackage{physics}
\usepackage{hyperref}
\usepackage{scalerel,stackengine}
\usepackage{fancyvrb}


\newcommand{\code}[1]{\texttt{#1}}
\newcommand{\shp}{$>>>$ }
\newcommand{\tab}{\hspace*{22pt}}

\usetheme{Madrid}

\title{Chapter 7: \\ Files and Modules}
\author{EEC 3150}
\date{}

\institute{Trevecca Nazarene University}
 


%%%%%%%%% BEGIN DOCUMENT %%%%%%%%%
%%%%%%%%% BEGIN DOCUMENT %%%%%%%%%
%%%%%%%%% BEGIN DOCUMENT %%%%%%%%%

\begin{document}

% get title page 
\frame{\titlepage}
 
%%%%%%%%%%
%%%%%%%%%%
\begin{frame}	
	\frametitle{Outline}
	\tableofcontents	
\end{frame}


%%%%%%%%%%
%%%%%%%%%%
\section{Files}

\begin{frame}	
	\center{\Huge{Files}}
\end{frame}

\begin{frame}	
	\frametitle{Files}
	
	\begin{block}{Intro to reading a file with Python:}
		\minipage{0.5\textwidth}
			\code{fname = "FirstFile.txt" \\
					\tab \\
					fhandle = open(fname, 'r')  \\
					\tab \\
					string = fhandle.read() \\
					print(string) \\
					\tab \\
					fhandle.close() 
			}
		\endminipage\hfill
		\minipage{0.5\textwidth}
			\begin{itemize}
				\item To read a file, use the \code{open()} function
				\item Two string arguments: 1) the filename, 2) whether to read (`r'), write (`w'), or append (`a') to the file
				\item Returns a \emph{file handle} (essentially a variable that points to the file)
				\item Since file handles are objects, the have their own methods (such as \code{read()}
				\item Close the file when finished
				\item Trying to read after closing throws an error
			\end{itemize}
		\endminipage\hfill
	\end{block}
	
\end{frame}

\begin{frame}	
	\frametitle{Files}
	
	\begin{block}{Reading files in a \code{with,as} block:}
		\minipage{0.6\textwidth}
			\begin{itemize}
				\item Closing a file is annoying and you can forget to do it
				\item For convenience, we can read files inside a \code{with,as} block
			\end{itemize}
			
			\vspace{20pt}
			\code{fname = "FirstFile.txt" \\
					with open(fname, 'r') as fhandle:
					\tab string = fhandle.read() \\
					\tab print(string)
			}
		\endminipage\hfill
		\minipage{0.4\textwidth}
			\begin{itemize}
				\item The file handle is only defined in the scope of the block
				\item The file is closed automatically after the block is finished
				\item Read the file in the block, save the contents to a variable, and use them later on
			\end{itemize}
		\endminipage\hfill
	\end{block}
	
\end{frame}

\begin{frame}	
	\frametitle{Files}
	
	\begin{block}{Some read methods:}
		\minipage{0.6\textwidth}
			\vspace{10pt}
			
			\code{with open(fname, 'r') as fhandle:
					\tab string = fhandle.read() \\
			}
			
			\vspace{10pt}
			\code{with open(fname, 'r') as fhandle:
					\tab lines = fhandle.readlines() \\
			}
		\endminipage\hfill
		\minipage{0.05\textwidth}
		\endminipage\hfill
		\minipage{0.35\textwidth}
			\code{fhandle.read()}
			\begin{itemize}
				\item Returns the entire file in one string
			\end{itemize}
			
			\vspace{10pt}
			\code{fhandle.readlines()}
			\begin{itemize}
				\item Returns a list of strings split on newlines
			\end{itemize}
		\endminipage\hfill
	\end{block}
	
\end{frame}

\begin{frame}	
	\frametitle{Files}
	
	\begin{block}{Reading a line at a time:}
		\minipage{0.6\textwidth}
			\vspace{30pt}
			
			\code{with open(fname, 'r') as fhandle:
					\tab for line in fhandle: \\
					\tab \tab print(line)
			}
			
			\vspace{30pt}
			\code{with open(fname, 'r') as fhandle:
					\tab line = fhandle.readline() \\
					\tab print(line) \\
					\tab line = fhandle.readline() \\
					\tab print(line)
			}
		\endminipage\hfill
		\minipage{0.05\textwidth}
		\endminipage\hfill
		\minipage{0.35\textwidth}
			With a for loop:
			\begin{itemize}
				\item We can read a line at a time with a for loop
			\end{itemize}
			
			\vspace{20pt}
			\code{fhandle.readline()}
			\begin{itemize}
				\item Returns the next line
			\end{itemize}
		\endminipage\hfill
	\end{block}
	
\end{frame}

\begin{frame}	
	\frametitle{Files}
	
	\begin{block}{Reading a line at a time:}
		\minipage{0.6\textwidth}
%			\vspace{30pt}		
			\code{with open(fname, 'r') as fhandle:
					\tab line = fhandle.readline() \\
					\tab while line: \\
					\tab \tab print(line) \\
					\tab \tab line = fhandle.readline()
			}
		\endminipage\hfill
		\minipage{0.05\textwidth}
		\endminipage\hfill
		\minipage{0.35\textwidth}
			With a while loop:
			\begin{itemize}
				\item We can read a line at a time with a while loop
				\item Read the first line
				\item Use the line as the loop condition
				\item Once \code{readline()} returns an empty string, the loop exits
			\end{itemize}
		\endminipage\hfill
	\end{block}
	
\end{frame}

\begin{frame}	
	\frametitle{Files}
	
	\begin{exampleblock}{Get a list of grades from a file:}
		Read the file Grades.txt and produce a list of integer grades.
		\begin{itemize}
			\item This file is simple and has one number per line
		\end{itemize}
	\end{exampleblock}
	
	\begin{exampleblock}{Get lists of years and temperatures from a file:}
		Read the file YearTemp.txt and produce lists of integer years and float temperatures.
		\begin{itemize}
			\item This file is more complicated
			\item It has a header that we must ignore
			\item Year and temp values are comma-separated (one pair per line)
		\end{itemize}
	\end{exampleblock}
	
\end{frame}




\begin{comment}


\begin{frame}%[fragile]
	\frametitle{Mutability}
	
	\begin{block}{Let's compare these codes:}
		\begin{itemize}
			\item These codes demonstrate defining a tuple/list \code{y} in terms of another \code{x} and then changing \code{x} afterwards
			\item In which codes does changing \code{x} change \code{y}? \\
					(\code{id()} returns an object's location in memory)
		\end{itemize}
		
		\vspace{10pt}
		\minipage{0.5\textwidth}
			\begin{footnotesize}
			\code{x = (1,) \\
			y = x \\
			print(y, x) \\
			print( id(y), id(x) ) \\
			\tab \\
			x = (1,2) \\
			\# x += (2,) \\
			\tab \\
			print(y, x) \\
			print( id(y), id(x) )
			}
			\end{footnotesize}
		\endminipage\hfill
		\minipage{0.5\textwidth}
			\begin{footnotesize}
			\code{x = [1] \\
			y = x \\
			print(y, x) \\
			print( id(y), id(x) ) \\
			\tab \\
			x = [1,2] \\
			\# x += [2] \\
			\# x.append(2) \\
			print(y, x) \\
			print( id(y), id(x) )
			}
			\end{footnotesize}
		\endminipage\hfill
	\end{block}
	
\end{frame}


\end{comment}




\end{document}




